\section{Effect of Evaluation Bias on Reported Performance}\label{sec:effect}
In this section, we provide two simple examples that illustrate how label ambiguity and temporal structure change the reported performance, and we provide intuition on why evaluation protocols differ in difficulty.

\subsection{Label Ambiguity}
Consider a feature vector $v$ that occurs 60 times with label $a$ and 40 times with label $b$ in a test set. A classifier assigns a single label to $v$, and therefore it classifies at most 60 of these 100 samples correctly. The contribution of $v$ to the accuracy is bounded by 0.6, independent of the learning algorithm and of the search effort. If such vectors cover a substantial share of the file, as in the traffic data, the same bound applies to the overall score.

On the raw data of device 9, the bound of Equation~\eqref{eq:ceiling} is 0.718. The best method reaches 0.711 (Section~\ref{sec:evaluation}), which is 99\% of the bound. In this regime, a ranking of frameworks reflects tie-breaking between labels of ambiguous vectors and not the ability to detect anomalies.

\textbf{Interpretation:} Ambiguity changes the meaning of a score, because the achievable accuracy depends on the file and not on the classifier. A high accuracy on a file with ambiguous vectors indicates that the classifier learned the majority labels, and a low accuracy does not necessarily indicate a classifier of low quality. Duplicated records enlarge this effect, because the majority label of an ambiguous vector is then repeated in both partitions.

\subsection{Temporal Structure}
Consider a recording that is divided into blocks of consecutive samples. Under a random protocol, a test sample at time $t$ has training neighbors at $t-1$ and $t+1$. If neighboring measurements are similar, a classifier can interpolate between neighbors and achieves a high score without learning the boundaries between operating conditions. Under a block protocol, the neighbors of a test block are absent from the training set, and the classifier must generalize to operating conditions that it has not observed.

A preliminary probe on MEMS illustrates the size of this effect for a single learning algorithm. A random forest with 100 trees and a maximum depth of 10 reaches an accuracy of 0.71 and a macro-F1 of 0.66 under a random 70/30 split. When whole blocks are held out in five-fold validation, the values fall to $0.54\pm0.12$ and $0.35\pm0.10$.

\textbf{Interpretation:} Both effects inflate the reported performance independent of the framework. A comparison of frameworks is therefore meaningful only after the protocol has been audited and fixed. We implement this requirement in the pipeline of the next section.
